%% ****** Start of file apstemplate.tex ****** %
%%
%%
%%   This file is part of the APS files in the REVTeX 4 distribution.
%%   Version 4.1r of REVTeX, August 2010
%%
%%
%%   Copyright (c) 2001, 2009, 2010 The American Physical Society.
%%
%%   See the REVTeX 4 README file for restrictions and more information.
%%
%
% This is a template for producing manuscripts for use with REVTEX 4.0
% Copy this file to another name and then work on that file.
% That way, you always have this originald template file to use.
%
% Group addresses by affiliation; use superscriptaddress for long
% author lists, or if there are many overlapping affiliations.
% For Phys. Rev. appearance, change preprint to twocolumn.
% Choose pra, prb, prc, prd, pre, prl, prstab, prstper, or rmp for journal
%  Add 'draft' option to mark overfull boxes with black boxes
%  Add 'showpacs' option to make PACS codes appear
%  Add 'showkeys' option to make keywords appear

\documentclass[preprint]{revtex4-1}
\usepackage{gnuplottex}
\usepackage{mathtools}
%\usepackage[usenames,dvipsnames]{xcolor}
\usepackage{float}
\usepackage{array}
\usepackage[english]{babel}
\usepackage[utf8]{inputenc}
\usepackage{amsmath}
\usepackage[table]{xcolor}
\usepackage{amsfonts}
\usepackage{graphicx}
\usepackage[colorlinks=true, linkcolor=blue, citecolor=dred]{hyperref}
\usepackage{algorithm}
\usepackage{algpseudocode}
\usepackage{bbold}
\usepackage{mathtools}
\usepackage{amssymb}
\usepackage{bm}
\usepackage{bbm}
\usepackage{tabularx, booktabs}
\usepackage{xspace}
\usepackage{listings}
\newcolumntype{Y}{>{\centering\arraybackslash}X}

\newcommand{\ket}[1]{\ensuremath{|#1\rangle}\xspace}
\newcommand{\bra}[1]{\ensuremath{\langle #1|}\xspace}
\newcommand{\psh}[2]{\ensuremath{\langle #1|#2\rangle}\xspace}
\newcommand{\etal}{{\it et al.}}
\newcommand\bbone{\ensuremath{\mathbbm{1}}}
\newcommand{\subsupi}[3]{#1_{\rm #2}^{\rm #3}}
\newcommand{\bfn}{\mathbf{n}}
\newcommand{\bfr}{\mathbf{r}}
\newcommand{\bfx}{\mathbf{x}}
\newcommand{\bfv}{\mathbf{v}}
\newcommand{\rmi}{{\rm i}}
\newcommand{\Hxc}{{\rm Hxc}}
\newcommand{\rmc}{{\rm c}}
\newcommand{\ddroit}{{\rm d}}
\newcommand{\undern}{\underline{n}}
\definecolor{dgreen}{rgb}{0,.5,0}
\definecolor{dblue}{rgb}{0,0,.5}
\definecolor{dred}{rgb}{0.5,0,.5}

\newcommand{\bruno}[1]{{\textcolor{blue}{Bruno: #1 }} }
%%%%%%%%%%%%%%%%%%%%%%%%%%%%%%%
\begin{document}

\title{Manual to generate intrinsic fragment and bonding orbitals within ROSE
}

\author{Bruno Senjean}
\email{bsenjean@gmail.com}
\author{Souloke Sen}
\author{Michal Repisky}
\author{Gerald Knizia}
\author{Lucas Visscher}

%%%%% Abstract %%%%%

\begin{abstract}
ROSE (Reduction of Orbital Space Extent) aims at constructing intrinsic fragment (IFOs)
and localized molecular orbitals (ILMOs).
It is mainly based on the construction of intrinsic atomic orbitals
(IAOs) and bonding orbitals (IBOs) from Ref.~\cite{knizia2013intrinsic}, but generalized to molecular fragments
and complex and quaternion algebra used in relativistic quantum chemistry as documented in Ref.~\cite{ROSE1}.
The idea is to define an intrinsic minimal basis set
of polarized orbitals which span the occupied space
exactly. The aforementioned orbitals are obtained by operating projections
between the molecular orbitals of the full molecule
and a minimal basis set of molecular orbitals obtained from
fragment (atomic or molecular) calculations.
A new feature introduced in 2023 and documented in  Ref.~\cite{ROSE2} is to apply the same procedure
also to localize additional "hard" virtuals. These will resemble a set of reference virtuals defined on the fragments,
in mathematically exactly the same way as the localized occupied orbitals resemble the occupied reference fragment orbitals.
\end{abstract}

\maketitle

\tableofcontents

\section{How to install}

\subsection{Dependencies}

ROSE is a standalone module 
interfaced with several quantum chemistry code.
For the moment, interfaces with 
Psi4~\cite{psi4}, PySCF~\cite{pyscf}, Gaussian~\cite{Gaussian09} (through formatted checkpoint files) and ADF~\cite{AMS19}
are implemented. An interface with ReSpect~\cite{respect}
is also on-going.

\subsection{Install ROSE}

Install ROSE with the commands specified in the README document at the top level directory of ROSE.
Note that the test folder and \verb|pip| installation is only useful if you want to try out the
python scripts/packages, which allow to run
all SCF calculations (for the full molecule and all fragments)
and ROSE in a single script for different interfaces.

\section{How to run}

ROSE can be run through python scripts, which is really convenient as
it performs everything for you in one single step: generation of the molecular and fragment
inputs (by running your favorite program...), generation of intrinsic fragments orbitals (IFOs) and intrinsic localized molecule orbitals (ILMOs) by executing ROSE,
and restart SCF calculations to check if the latters are indeed correct.
But obviously, if you already have all the inputs, no need to re-do the SCF calculations again !
You can just set all your input files in a given folder, write an \verb|INPUT_GENIBO| file
and execute ROSE to get the resulting IFOs and ILMOs.
The format for the input files are detailed in the following.
Note that the provided scripts are sufficient for standard applications of ROSE, but not if one wants to have more flexibility.
{\bf Therefore, we strongly advise the user to first look and test the
different examples provided in the ROSE directory. If one is willing to have more flexibility, then reading this manual will be useful.}

\subsection{Warnings}

All calculations have to be performed without symmetry (symmetry \verb|c1|) and
with uncontracted basis sets and
cartesian functions.

\subsection{Guide through the code}

ROSE is a FORTRAN program. Its main program is \verb|generate_ibo.F90| and
it uses several modules, all in the \verb|src| folder:
\begin{itemize}
\item \verb|linear_algebra.F90| to do linear algebra operations for real, complex and quaternion matrices.
\item \verb|rose_ao.F90| and \verb|rose_mo.F90| to define the atomic and molecular orbital data types. They define all the necessary information to manipulate molecular orbitals.
\item \verb|utils.F90| contains normalization factors for different interfaces as well as different useful functions.
\item \verb|import_fragment.F90| define atomic or molecular fragments orbitals (number of core and valence orbitals) and contains several functions useful functions.
\item \verb|dirac_interface.F90|, \verb|hdf5_interface.F90| and
\verb|fchk_interface.F90| handle the reading and writing of basis set and orbital informations within the different interfaces. The default is to use
the hdf5 interface, for which we have python scripts to populate it with the information needed by pyscf or psi4. At the moment of writing, the DIRAC hdf5 interface does still
have some minor differences compared to these non-relativistic codes, but this is planned to be unified soon. 
\item \verb|make_iao.F90| and \verb|make_ibo.F90| are subroutines to compute intrinsic fragment orbitals (IFOs) and intrinsic localized molecular orbitals (ILMOs), respectively. Also contains routines for re-canonicalization.
\item \verb|properties.F90| is a module to compute properties like the percentage of delocalization and partial charges.
\item \verb|check.F90| check the orthonormalization of the resulting IFOs and ILMOs basis.
\item \verb|rose_input_data.F90| Module to hold the input parameters of ROSE. These can be set by the user from input.
\end{itemize}

Several input files are needed before executing \verb|genibo.x| (the 
executable of \verb|generate_ibo.F90|). The only input file that is common to any
interface is \verb|INPUT_GENIBO| and contains the following keywords that can be given in any order inside the input block starting with \verb|**ROSE| and ending
with \verb|*END OF INPUT|. Keywords themselves always start with a period and may be followed by a line in which a numerical value (real or integer) is given.

\begin{itemize}
\item{Keywords defining the molecule to be fragmented and the reference orbitals (core, valence, virtual).}
\begin{lstlisting}[language=bash]
.CHARGE
.NFRAGMENTS
.FILE_FORMAT
.FRAG_FORMAT
.UNRESTRICTED
.SPINORS
.FRAG_VALENCE
.FRAG_CORE
.ADDITIONAL_VIRTUALS_CUTOFF
.FRAG_THRESHOLD
.RANKTOL
.TIETOL
\end{lstlisting}
\item{Keywords defining the algorithm to be used.}
\begin{lstlisting}[language=bash]
.VERSION
.EXPONENT
.INCLUDE_CORE
.FRAG_BIAS
.PREPROJ
\end{lstlisting}
\item{Keywords controlling optional additional procedures.}
\begin{lstlisting}[language=bash]
.AVAS_THRESHOLD
.FRAG_AVAS
.TEST
\end{lstlisting}
where 
\begin{itemize}
\item \verb|CHARGE| specifies the charge of the whole molecule. Default is 0
\item \verb|NFRAGMENTS| specifies number of molecular fragments. For atomic fragments set to 0. Default is 0
\item \verb|FILE_FORMAT| can be \verb|DIRAC|, \verb|hdf5|, or \verb|adf| depending on the format of the MO coefficient 
file of the full molecule (basis $\mathcal{B}_1$).
\item \verb|FRAG_FORMAT| can be \verb|DIRAC|, \verb|hdf5|, or \verb|adf| 
with as default format for the fragment (basis $\mathcal{B}_2$) the same format as for the full molecule.
\item \verb|UNRESTRICTED| specifies unrestricted type calculation, default is restricted.
\item \verb|SPINORS| specifies to use spinors for 
relativistic calculations. Default is  to use spatial orbitals (alpha and beta spinors share the same spatial orbital)
\item \verb|FRAG_VALENCE| specifies the number of valence orbitals for each fragment. If not specified, by default the code generates the valence space for each fragment as the sum of valence orbitals of individual atoms within each fragment. Note that this block needs to be repeated for every fragment the user wishes to specify. The structure of the block is as follows: \\
\verb|FRAG_VALENCE| \\
\textit{fragment \#}\\
\textit{\# of valence orbitals}\\
\item \verb|FRAG_CORE| specifies the number of core orbitals for each fragment. If not specified, by default the code generates the core space for each fragment as the sum of core orbitals of individual atoms within each fragment. It has the same block structure as \verb|FRAG_VALENCE|.
\item \verb|ADDITIONAL_VIRTUALS_CUTOFF| specifies the energy threshold (in hartrees) to determine the number of additional hard virtuals to localize for the supermolecule. Default value of this cutoff value is 2.0 hartrees.
\item \verb|FRAG_THRESHOLD| specifies energy threshold (in hartrees) to determine number of reference hard virtuals to use for the fragments (needed to localize the additional hard virtuals). The total number of reference virtuals should be larger 
than the number of virtual MOs to localize (similar to having more reference valence orbitals than occupied valence orbitals to localize), so it is advisable to use a relatively high value for this parameter.  Default is 10.0 hartrees.
A generous value used to be dangerous, because a reference set larger than the complementary space of the supermolecule stopped the calculation. The reference virtuals are now reduced to a subset that stays linearly independent after projection on that space (see \verb|RANKTOL|), so the number selected here is an upper bound rather than a quantity that has to be tuned by hand.
\item \verb|RANKTOL| specifies the threshold below which a reference orbital is considered to add nothing once it has been projected on the space it is meant to describe. It governs two things: the reduction of the reference virtuals to a linearly independent subset, and, under \verb|PREPROJ|, the check that the valence reference set of each fragment keeps its rank. Default is 1.0D-6, which discards numerically dependent directions while leaving everything else alone; raising it to 1.0D-4 or 1.0D-3 trades a little of the span for a better conditioned reference set, while lowering it to 1.0D-8 retains more of the span at the price of a less well conditioned one.
\item \verb|TIETOL| specifies the relative width of the group of candidates treated as tied during that reduction. Symmetry-equivalent fragments yield exactly equal residuals, so taking them one at a time would decide between them on rounding noise and give equivalent fragments different numbers of reference orbitals. All candidates within this relative distance of the best one are therefore taken together. Default is 1.0D-4, which is loose enough to recognise fragments that were converged independently. It is a tie tolerance and is meant to stay small: a value approaching 1 makes every remaining candidate a tie and degrades the selection to the order in which the reference orbitals happen to be supplied.

\item \verb|VERSION| specifies version of the IAO generation. can be either \verb|Stndrd_2013|, \verb|Simple_2013| or
\verb|Simple_2014| and corresponds to the version of the transformation to go from the
RFOs to the IFOs.
Those versions were designed by Knizia in his original Python program and may lead to small differences in the functions that are generated. More information
about the precise procedures can be found in the SI of our first ROSE paper\cite{ROSE1}.
\item \verb|EXPONENT| is the value of the exponent (for the default Pipek-Mezey type of localization, the exponent is equal to 2. A higher value of 3 or 4 is sometimes useful.).
\item \verb|INCLUDE_CORE| specifies to include core orbitals in the generation of IAOs and localization of the occupied orbitals. Default is false, which means the first IAOs are actually the (delocalized) core orbitals of the full molecule, same of IBOs.
Note that this option is overwritten by \verb|FRAG_CORE| and \verb|FRAG_VALENCE|
\item \verb|PREPROJ| specifies to project the reference fragment orbitals on the orbital space of the supermolecule before the IFOs are constructed, so that the reference space is contained in the molecular one by construction. This removes whatever part of the reference orbitals the molecular basis cannot carry, which matters when the fragments come from a library computed in a larger basis than the molecule itself. It also makes the variants of the IFO construction selected by \verb|VERSION| coincide, since the projector on the molecular space then acts on the reference set as the identity. A fragment whose own valence reference set loses rank under the projection stops the calculation, as its reference basis cannot be represented in the molecular one. Default is false.
\item \verb|FRAG_BIAS| specifies the bias applied to each fragment when assigning the localized orbitals to a particular fragment. This allows to identify target fragments to which covalent bonds that connect two or more fragments should be preferably assigned.
A default bias of 0.1 is applied to fragment 1 which is sufficient to assign non-polar covalent bonds to this fragment. Follows the same structure as \verb|FRAG_VALENCE|.

\item \verb|FRAG_AVAS| specifies the fragments and orbitals to be used for projecting and definition of the automated valence active space ~\cite{sayfutyarova2017automated}. Default is an empty list.
If you want to project on orbitals 8 through 10 of fragment 2, the input is first a line with a \verb|2|, followed by a \verb|3| (the number of orbitals), followed by the numbers \verb|8|, \verb|9|, and \verb|10| in free format.
This keyword can be repeated to allow projection on more than one set of fragment orbitals. The sets are concatenated so that only one active space will be produced.
\item \verb|TEST| specifies to perform some tests on the correctness of the localized occupied orbitals: orthonormalization of matrices and deviation in the number of electrons are checked, etc. 
Default is false.
\end{itemize}
Note while all keywords with a specified default value are optional,  specification of the file format is mandatory.
{\it Example}. You are generating the IFOs of the water molecule using psi4 and atomic fragments. A minimal input file is given below:
\begin{lstlisting}[language=bash]
**ROSE
.FILE_FORMAT
hdf5
*END OF INPUT
\end{lstlisting}

Once all the input files are all in the same folder, the only remaining step is
to execute ROSE:
\begin{lstlisting}[language=bash]
/path_to_rose/bin/genibo.x
\end{lstlisting}

Thie code will print information on the localization procedure, and will generate
IFO and IBO files (with appropriate extension defined in the input).
It is of course convenient to have a simple way of plotting the orbitals. This can be done in \verb|AVOGADRO| 
for instance with FCHK files. ROSE will provides a python notebook to convert from hdf5 to fchk format (work in progress). 
Note that it only works for real spatial orbitals as visualization of complex spinors is non-trivial.
In the following, we show in more details what are the required input files, and how they can be generated by a simple python interface.

\subsection{Generating the input files}
For the provided scripts, several input files are required:
\begin{itemize}
\item geometry files with extension \verb|.xyz| (in cartesian coordinates)
\begin{enumerate}
\item MOLECULE.XYZ (the geometry file of the full molecule)
\item If \verb|NFRAGMENTS=0|, one file for each different type of atom, all centered 
at position (0,0,0). The files are named as \verb|"{:03d}.xyz".format(Z)| where \verb|Z|
is the atomic number of the atom. For instance, hydrogen will be \verb|001.xyz| while
carbon is \verb|006.xyz|.
\item  If \verb|NFRAGMENTS != 0|, one file for each different fragments, with the
same geometry as in the full molecule. The files are named as 
\verb|"frag{:d}.xyz".format(i)| where \verb|i|
is the label of the fragment. For instance, for the case of water-ammonia, one can define two 
molecular fragments (water and ammonia) called \verb|frag0.xyz| and \verb|frag1.xyz|.
\end{enumerate}
\item MO coefficient files
\begin{enumerate}
\item In DIRAC, those are \verb|h5| files corresponding to DIRAC23 and later. The DFCOEF files have to be renamed in the same way as 
geometry files, like \verb|MOLECULE.h5|, \verb|001.h5| or \verb|frag0.h5|.
\item For formatted checkpoint files with extension \verb|.fchk|,
it might
be necessary to first truncate it to keep only the necessary informations to be read in ROSE.
This can be done thanks to the python script in \verb|/rose/python_scripts/|, called
\verb|python_transform_fchk.py|.
\item A python package can interface Psi4 with ROSE by creating a \verb|.h5| extension which looks like the \verb|.h5| one of DIRAC. The resulting localized orbitals are also written in this format and can easily be
fed back into via a Python script. A similar interface has been implemented for the \verb|.pyscf| extension.
\item in ADF, the rkf  (extension \verb|.rkf|) files are used.
\end{enumerate}
\item Finally, the \verb|INPUT_GENIBO| file which has already been described.
\end{itemize}

The ADF interface is quite different from the other ones.
This is caused by the use of Slater-type Orbitals (STOs)
instead of Gaussian-type Orbitals (GTOs) by ADF.
Thus, the InteRest library~\cite{interest} of ROSE to compute overlap matrices in the AO basis
cannot be used; \verb|slater_overlap.F90| computes them instead, directly from the
Slater basis, so that nothing has to be taken from ADF for this purpose.

ROSE itself no longer reads or writes ADF's binary KF files. An ADF calculation
is converted to the HDF5 format of Sec.~\ref{sec:kfbridge} and read from there,
which is why a standalone ROSE needs neither a KF library nor an AMS
installation. The workflow below is still how the ADF calculations themselves are
produced. Before executing
\begin{lstlisting}[language=bash]
/path_to_adf/executable/plams AMS.py
\end{lstlisting}
one has to create the \verb|INPUT_ADF| file as follows
\begin{lstlisting}[language=bash]
/path_to_the_workdir/directory/
RESTRICTED
OPENSHELL
BASIS1
BASIS2
CHARGE
MULTIPLICITY
NUMBER_OF_FRAGMENTS
/path_to_frag0.xyz nMO CHARGE MULT
/path_to_frag1.xyz nMO CHARGE MULT
RESTART
\end{lstlisting}
\end{itemize}
where
\begin{itemize}
\item \verb|/path_to_the_workdir_directory/| contains \verb|MOLECULE.xyz| (note the lower 
case for \verb|xyz| here)
\item \verb|RESTRICTED| is either True or False
\item \verb|OPENSHELL| is either True or False
\item \verb|BASIS1| is the string name of the basis of the full molecule, like \verb|SZ| or 
\verb|DZP| (see ADF documentation)
\item \verb|BASIS2| is the string name of the basis of the fragments
\item \verb|CHARGE| is the charge of the full molecule
\item \verb|MULT| is the spin multiplicity of the full molecule
\item \verb|NUMBER_OF_FRAGMENTS| is the number of fragments
\item And then one line for each fragment, containing the path to its \verb|.xyz| file, 
the number of MOs \verb|nMO| (minimal truncated set which will sum up to the final number of IFOs),
the charge and the spin multiplicity, separated by a space.
\end{itemize}

If everything is set up correctly and if the required \verb|xyz| files are in the proper folder,
running \verb|plams AMS.py| will perform all the SCF calculations and generate the
necessary rkf files. These are then converted as described in
Sec.~\ref{sec:kfbridge}, and ROSE reads the converted files.
Examples of python scripts for running the entire workflow starting from the SCF calculations, generation of the  \verb|INPUT_GENIBO|, execution of  \verb|genibo.x| to the post-SCF calculations can be found in the \verb|examples| directory for each of the interfaces with DIRAC,ADF,PSI4 and PYSCF.

\subsection{Moving data between ADF and ROSE: \texttt{rose\_kf\_bridge.py}}
\label{sec:kfbridge}

ADF keeps its data on binary KF files, while ROSE reads and writes HDF5. The two
are connected by \verb|python_scripts/rose_kf_bridge.py|, a developer tool run
under \verb|amspython|. It is deliberately kept outside the build: neither the
ROSE executable nor the test suite ever calls it, so a standalone ROSE needs no
KF library and no AMS installation. Only this script does, and only when it is
used.

\paragraph{Forward: rkf to HDF5.}
\begin{lstlisting}[language=bash]
amspython rose_kf_bridge.py forward adf.rkf \
          -o CHECKPOINT.h5 --t15 TAPE15
\end{lstlisting}
reads the Slater basis, the MO coefficients and the geometry from an ADF
calculation and writes them in the format \verb|read_from_hdf5| expects. ADF
stores Cartesian-power Slater functions,
$\chi = N\,x^{k_x}y^{k_y}z^{k_z}r^{k_r}e^{-\alpha r}$ with $k_x+k_y+k_z = l$ and
$k_r = n-1-l$, so the basis is written with \verb|angular = 1| and
\verb|basis_type = 2|, with one extra array for the principal quantum number per
primitive. The MO coefficients are permuted out of ADF's own AO order into the
order of the shell list, using \verb|A%npart|.

With \verb|--t15| the script also stores the AO overlap matrix that ADF itself
computed, read from \verb|Matrices%Smat|. ROSE compares its Slater overlaps against it whenever that
dataset is present, which checks the AO ordering, the normalisation, the angular
convention and the overlap routine in a single number. Because the reference
travels inside the converted file, the check keeps working on machines that have
never had ADF installed.

The converted test inputs stored under the \verb|input| directory of each of
the two ADF tests were produced this way: one call for the supermolecule and
one per fragment, the fragment files being named by atomic number for atomic
fragments, and \verb|frag1.h5|, \verb|frag2.h5| and so on for molecular ones.

\paragraph{Backward: HDF5 to rkf.}
\begin{lstlisting}[language=bash]
amspython rose_kf_bridge.py backward ribo.h5 \
          -o Molecule_B1_MO1_ribo.rkf \
          --template Molecule_B1_MO1.rkf
\end{lstlisting}
copies the template rkf and writes a \verb|Rose Data| section into the copy,
holding the localized orbitals in ADF's AO order together with the orbital
energies, the occupations, the counts per recanonisation group and the RIBO
Fock matrix. This is what ADF's TDDFT reads.

It should be understood what that file is. \textbf{The result of a ROSE
calculation is \texttt{ribo.h5}}; the rkf written here is an \emph{intermediary},
produced so that an ADF calculation can be continued in the orbital basis ROSE
has constructed. It is a re-encoding of \verb|ribo.h5| for one consumer, it
carries no information that \verb|ribo.h5| does not, and it should not be
archived or cited as the record of the calculation. Regenerate it from
\verb|ribo.h5| whenever ADF needs it.

\paragraph{When something is unclear.} \verb|inspect| lists the sections and
variables of a KF file that the conversion depends on, and \verb|probe| prints
the signatures of the \verb|KFFile| bindings and tries each way of opening a
file. Both are read-only.

\section{Contact}

If you experience any trouble, please contact one of the authors \verb|bsenjean@gmail.com| or \verb|l.visscher@vu.nl|

\bibliographystyle{apsrev4-1}
\bibliography{biblio}

\end{document}
